\documentclass[11pt]{article}
\usepackage[a4paper,margin=26mm]{geometry}
\usepackage[T1]{fontenc}
\usepackage{lmodern}
\usepackage{amsmath,amssymb,amsthm}
\usepackage{microtype}
\usepackage{array}
\usepackage{xurl}
\usepackage{hyperref}
\hypersetup{colorlinks=true,linkcolor=blue,urlcolor=blue,citecolor=blue,
 pdftitle={Two sufficient conditions for a tree to be determined by its chromatic symmetric function},
 pdfauthor={Horace Dong}}
\newtheorem{theorem}{Theorem}
\renewcommand{\thetheorem}{\Alph{theorem}}
\newtheorem{lemma}{Lemma}
\newtheorem{proposition}[lemma]{Proposition}
\newtheorem{corollary}[lemma]{Corollary}
\theoremstyle{definition}
\newtheorem{definition}[lemma]{Definition}
\newtheorem{example}[lemma]{Example}
\theoremstyle{remark}
\newtheorem{remark}[lemma]{Remark}
\newcommand{\repoversion}{1.27.0}
\setlength{\emergencystretch}{1.5em}
\newcommand{\N}{\mathbb N}

\title{Two sufficient conditions for a tree to be determined\\ by its chromatic symmetric function}
\author{Horace Dong\\ \small Independent researcher\\
\small ORCID: \href{https://orcid.org/0009-0002-7554-7548}{0009-0002-7554-7548}}
\date{28 November 2026}
\begin{document}
\sloppy
\maketitle

\begin{abstract}
Stanley asked in 1995 whether two non-isomorphic finite trees can have the same
chromatic symmetric function~$X_T$; the question is open. For a tree $T$, let
its core be the subtree of non-leaf vertices, and give each core vertex the
weight $1$ plus its number of leaf neighbours. We prove two sufficient
conditions, both stated only for $T$: every tree with the same $X$ is isomorphic
to $T$ if (A) every core vertex has a leaf neighbour, every leaf of the core has
a weight that occurs only once in the core, and any two sets of core vertices
with the same total weight contain the same leaves of the core; or if (B) every
core vertex has a leaf neighbour, no core weight is the sum of the weights of
two or more other core vertices, and at most one leaf of the core has a weight
that occurs elsewhere in the core. Both conditions allow internal core weights
to coincide, also at adjacent vertices, unlike the recent conditions of Zeng
and of Awan, which require unique weights at internal core vertices. The
reconstruction step for (A) rebuilds the core from data indexed by the core
leaves with Buneman's classical theorem, a device used earlier, for a different
class allowing one connected class of repeated weights, in a public preprint of
Atwell. The ingredients that $X_T$
determines the number of leaves, properness and the multiset of core weights are
known. The two classes are thin: family~(A) has $O(\log n)$ core leaves and
family~(B) fewer than $\sqrt{2n+2}$; every member with at most $29$ vertices is
already covered by the computations of Heil and Ji, every member of diameter six
is covered by two public results posted later on the day this work was completed
(one of them with a Lean proof; both unrefereed, and we have not checked them),
and the general question remains open. Both theorems are formally verified in Lean~4, with $X_T$ defined
through proper colourings. The proofs and the formalization were produced by AI
agents working under the author's direction; their role is described at the end
of the paper.
\end{abstract}

\section{Introduction}\label{sec:intro}

For a finite simple graph $G$ with vertex set $V$, Stanley's \emph{chromatic
symmetric function}~\cite{stanley} is
\[
X_G=\sum_{\kappa}\ \prod_{v\in V}x_{\kappa(v)},
\]
the sum over all proper colourings $\kappa:V\to\{1,2,3,\dots\}$ (adjacent
vertices receive different colours), in commuting variables $x_1,x_2,\dots$.
Stanley~\cite{stanley} asked whether non-isomorphic trees always have different
chromatic symmetric functions; this is often called Stanley's tree isomorphism
conjecture, and it is open. It has been verified by computer for all trees with
at most $29$ vertices (Heil and Ji~\cite[Theorem~4.5]{hj}), and proved for several
classes: spiders and some caterpillars (Martin, Morin and
Wagner~\cite{mmw}), proper caterpillars (Aliste-Prieto and Zamora~\cite{apz}),
all caterpillars (Loebl and Sereni~\cite[Theorem~2]{ls}), proper trees of
diameter at most~$5$ (Aliste-Prieto, de Mier, Orellana and
Zamora~\cite[Corollary~7.14]{apdmoz}), and all trees of diameter at most~$5$
(Gonzalez, Orellana and Tomba~\cite[Theorem~5.20, Proposition~5.24 and
Theorem~5.27]{got}). Orellana and
Tomba~\cite[Corollary~4.12]{ot} showed that it suffices to prove the conjecture
for proper trees all of whose core weights (defined below) are equal to a fixed
$k\ge2$.

\paragraph{Cores and weights.} Let $T$ be a finite tree with at least three
vertices. Its leaves are the vertices of degree~$1$; its \emph{core} $K(T)$ is
the subtree induced by the other vertices, and the \emph{leaves of the core} are
the vertices of degree~$1$ in $K(T)$ (they are not leaves of $T$). Every leaf of
$T$ has exactly one neighbour, which lies in the core. The \emph{weight} of a
core vertex $v$ is
\[
w(v)=1+\#\{\text{leaves of $T$ adjacent to $v$}\},
\]
so that $w(v)$ is the number of vertices in the star formed by $v$ and its leaf
neighbours, and $\sum_{v\in K(T)}w(v)=n$, the number of vertices of $T$. We write
$w(P)=\sum_{v\in P}w(v)$ for $P\subseteq K(T)$ and $W(T)$ for the multiset of
core weights. The tree $T$ is \emph{proper} if every core vertex has a leaf
neighbour, that is, $w(v)\ge2$ for all core vertices. Conversely, $T$ is
determined up to isomorphism by its weighted core: attach $w(v)-1$ new leaves to
each core vertex~$v$.

\paragraph{The results.} The hypotheses are imposed only on $T$; the tree $T'$
is arbitrary.

\begin{theorem}[anchor family]\label{thm:A}
Let $T$ be a finite tree whose core $K$ has at least two vertices, and assume:
\begin{itemize}
\item[(A1)] $T$ is proper;
\item[(A2)] every leaf $a$ of the core has a weight that no other core vertex
has: $w(v)=w(a)$ with $v\in K$ implies $v=a$;
\item[(A3)] if $P,Q\subseteq K$ and $w(P)=w(Q)$, then $P$ and $Q$ contain the
same leaves of the core.
\end{itemize}
Then every finite tree $T'$ with $X_{T'}=X_T$ is isomorphic to $T$.
\end{theorem}

\begin{theorem}[peeling family]\label{thm:B}
Let $T$ be a finite tree whose core $K$ has at least two vertices, and assume:
\begin{itemize}
\item[(B1)] $T$ is proper;
\item[(B2)] no core weight is a sum of core weights over a set of at least two
vertices: $w(v)\ne w(P)$ for all $v\in K$ and all $P\subseteq K$ with
$|P|\ge2$;
\item[(B3)] at most one leaf of the core has a weight that also occurs at
another core vertex.
\end{itemize}
Then every finite tree $T'$ with $X_{T'}=X_T$ is isomorphic to $T$.
\end{theorem}

\noindent The core leaves act as labels (``anchors'') in Theorem~\ref{thm:A}:
by~(A3), the total weight of a set of core vertices determines which core leaves
it contains. Condition~(B2) holds, for example, when all core weights lie in an
interval $[b,2b)$. Neither theorem requires the internal core weights to be
distinct: they may coincide, at adjacent vertices and with any multiplicity. Both theorems are
formally verified in Lean~4 (Section~\ref{sec:lean}), with $X$ defined through
proper colourings; the examples and estimates of Section~\ref{sec:examples} are
not part of the Lean statements.

\paragraph{Prior work and what is new.} The facts about $X_T$ used in the proofs
are known. That $X_T$ determines the number of vertices and the number of
leaves is in Martin, Morin and Wagner~\cite[Proposition~2(i) and~(v)]{mmw} (they also
recover the degree sequence~\cite[Corollary~5]{mmw}). That $X_T$ determines
whether $T$ is proper and, for proper $T$, the multiset $W(T)$ of core weights
follows from Gonzalez, Orellana and Tomba~\cite[Theorem~4.15 and
Corollary~4.20]{got}: the leading partition of $X_T$ in the star basis is the
partition of the orders of the leaf components, and it has no part~$1$ exactly
when $T$ is proper. That $X_T$ of a proper tree is computed from its weighted
core is in~\cite[Theorem~5.10 and Corollary~5.11]{apdmoz} (see
also~\cite[Proposition~3.8]{ot}), and coefficients that come from spanning
forests with a single non-singleton component are used in~\cite[Lemma~7.4]{apdmoz}.
The reconstruction step of Theorem~\ref{thm:A} is a version of the classical
fact that a tree is determined by its splits (Buneman~\cite{buneman}), with the
masses of the splits ordering the vertices of degree~$2$ in the core. This
method is not new: Atwell~\cite[Theorem~2.1 and Lemma~5.1]{atwell1} (a public
preprint of 12 November 2026, unrefereed and not formally verified) recovers from $X_T$ the distances between the tagged vertices of a connected
subtree $C$ of the core on which one weight repeats, among them all leaves of $C$,
and rebuilds $C$ from them by Buneman's theorem on additive tree metrics; Theorem~\ref{thm:A} recovers instead the splits of the
core labelled by its leaves and rebuilds the core from them.

Among the closest sufficient conditions are those of Zeng~\cite[Theorem~1.1]{zeng}: a
proper tree is determined by $X_T$ if every core vertex that is not a leaf of the
core has a weight occurring nowhere else in the core, while repetitions among core
leaves are unrestricted (``multiplicity-isolated'' cores); Awan~\cite[Theorem~1]{awan}
shows that for the same trees the weaker invariant
$\partial X_T/\partial p_1|_{p_1=0}$ already distinguishes them among proper
trees. Theorems~\ref{thm:A} and~\ref{thm:B} put the uniqueness on the core
leaves instead and leave the internal weights free, at the price of the
arithmetic conditions (A3) and (B2). Atwell's criterion~\cite[Theorem~2.1]{atwell1}
also allows repeated internal weights at adjacent vertices: exactly one weight $a$
may repeat, its vertices must form a connected subtree $C$ of the core, all
other weights are unique, every component of the rest of the core is attached to
$C$ by one edge with an isolated mass, every such peripheral mass is smaller
than~$a$, and every leaf of $C$ carries such an attachment. His class and our two
families are incomparable: his $69$-vertex example~\cite[Section~6]{atwell1}
satisfies (A1)--(A3) and (B1)--(B3) and has diameter~$7$; his class contains trees
in neither family (for example a star of four vertices of weight $20$, his $C$, with single
attached vertices of weights $2$, $3$ and~$5$, since $5=2+3$); and our examples of diameter~$7$ violate his condition
that the peripheral masses be smaller than~$a$ (Remark~\ref{rem:public}). An
earlier preprint of Atwell~\cite{atwell2} treats proper trees whose weights are
all congruent modulo some $m\ge2$ and whose repeated weights sit at pairwise
non-adjacent vertices, against proper competitors. If the internal weights of a tree in
family~(A) are pairwise distinct, then all its core weights are distinct by
(A2), and the tree is covered by~\cite{zeng,awan}; so family~(A) adds only trees
in which some internal core weight repeats. A tree that lies in family~(B) and
in Zeng's class has all core weights distinct (a repeated weight would have to
sit at two core leaves), so it is covered by~\cite[Corollary~5.15]{got}, which
treats proper trees with pairwise distinct core weights. The two families are
incomparable (Example~\ref{ex:incomp}). Neither contains the trees of the
reduction of Orellana and Tomba~\cite[Corollary~4.12]{ot}: if all core weights
are equal and the core has at least two vertices, then every core leaf violates
(A2) and at least two core leaves violate~(B3).

Both families are thin (Section~\ref{sec:examples}): a tree in family~(A) with
$n$ vertices has at most $\log_2(n+1)$ core leaves, and a tree in family~(B)
fewer than $\sqrt{2n+2}$. Every member with at most $29$ vertices is already
covered by the computation of Heil and Ji~\cite{hj}. In an exhaustive search
over cores with at most $8$ vertices, made by the final reviewer of
Theorem~\ref{thm:A} (see the end of the paper), the smallest member of
family~(A) that is not a caterpillar or a spider and has diameter at least~$6$
has $33$ vertices; for family~(B), the final reviewer of Theorem~\ref{thm:B}
found members with $30$ and $36$ vertices that lie outside all published and
public criteria it checked on 27 November (Example~\ref{ex:small}). These
examples have diameter~$6$. Two public results posted later that day state that
$X_T$ determines every such tree: a repository of K.~Feng~\cite{kyf} (created at
06:11~UTC, release~1.0.0 at 06:33~UTC, release~1.1.0 at 21:16~UTC) states, with a
manuscript and a Lean~4 proof that we have not compiled, that $X_T$ determines
every tree of diameter six among all trees; and a written proof without formal
verification~\cite[pull request~1081]{sc}, merged at 18:17~UTC, states this for
every proper tree of diameter six (its initial step, stated there for proper
trees of every diameter, recovers from $X_T$ the numbers of sets of core vertices
with given cardinality and weight and given numbers of core edges inside the set
and inside its complement). Both are unrefereed, and we have not checked them;
they cover all members of both families of diameter~$6$. Both families
have members of every larger diameter. Among the members that are not
caterpillars, the smallest of diameter at least~$7$ have $41$ vertices in
family~(A) and, outside Zeng's class, $21$ vertices in family~(B); family~(B)
also has members with $30$ vertices and diameter~$7$. The members with $41$ and
$30$ vertices given in Examples~\ref{ex:A} and~\ref{ex:small} satisfy none of the
explicit public criteria we checked (Remark~\ref{rem:public}), whereas other
members of diameter~$7$, such as Atwell's $69$-vertex example, are covered
by~\cite{atwell1}; the density and genericity results of~\cite{sc} do not decide
individual trees and were not used. The coloured reconstruction
lemma used for Theorem~\ref{thm:B} (Lemma~\ref{lem:colour}) is elementary; we
did not find it stated, but it may well be known.

Several public AI-assisted projects work on the same question; their results are
written, unrefereed and, unless stated otherwise, not formally verified. Working
notes in the pull request~\cite{aur} state sufficient conditions on repeated
weights, some of them with a similar device of distinct subset sums. The pull
requests~\cite{sc} of another project, merged on 27 November 2026, state
conditions based on long subdivisions, balanced edges, leaf-rich centroids, one
or two branch vertices with long arms, and arbitrary cores with long exterior
arms. One of them (pull request~1029, merged at 14:43~UTC, after the Lean proof of
Theorem~\ref{thm:A} was completed at 03:18~UTC) reconstructs arbitrary tree cores
carrying long exterior arms whose orders have separated subset sums; it turns cut
sizes into labelled splits by the same device as the reconstruction step of
Theorem~\ref{thm:A}, for a different class of trees: there every arm at a hub
has at least three vertices, so the hubs carry no pendant leaves, whereas every
vertex of degree at least three of a proper tree has a leaf neighbour.
Another (pull request~1081, merged at 18:17~UTC) is the diameter-six result
mentioned above, as is the repository~\cite{kyf}, which states the result for all trees of
diameter six with a Lean proof that we have not compiled. A preprint of
Carptopus~\cite{carp} treats trees with exactly three branch vertices. A third
condition, concerning trees with many leaves at a centroid, was studied in the
same work and frozen as a Lean statement at the same time, but a public draft
merged on 27 November 2026 at 04:28~UTC~\cite[pull request~977]{sc} covers it
with a weaker hypothesis; it is not part of this paper.

So the new content, as far as the searches described at the end of the paper
found, is the two sufficient conditions themselves, with complete proofs and a
Lean verification from the colouring definition of~$X$; every member they add
beyond earlier and public results has diameter at least~$7$ and more than $29$
vertices, and not all such members are new (Atwell's $69$-vertex example is
one). The reconstruction method of Theorem~\ref{thm:A} is not new
either~\cite{atwell1}. They are modest results: the arguments use only data from $X_T$
that is known to be accessible, the classes are thin, and they say nothing about
the reduced case of~\cite{ot} or about the general question.

\section{What the chromatic symmetric function determines}\label{sec:known}

For a tree $T$ and a partition $\lambda$ of $n=|V(T)|$, let $N_T(\lambda)$ be the
number of partitions of $V(T)$ into blocks that induce connected subgraphs and
have sizes~$\lambda$. In a tree, such partitions correspond bijectively to sets
of edges (the edges inside the blocks), and the block sizes are the orders of
the components of the spanning forest formed by these edges.

\begin{lemma}[Stanley]\label{lem:stanley}
$X_T$ determines $n$ and $N_T(\lambda)$ for every partition $\lambda$ of~$n$.
\end{lemma}

\begin{proof}
By~\cite[Theorem~2.5]{stanley}, $X_T=\sum_{S\subseteq E}(-1)^{|S|}p_{\lambda(S)}$,
where $p_\lambda$ are the power sums and $\lambda(S)$ is the partition of the
component orders of $(V,S)$. A forest with $\ell(\lambda)$ components has
$n-\ell(\lambda)$ edges, so the coefficient of $p_\lambda$ is
$(-1)^{n-\ell(\lambda)}N_T(\lambda)$, and the power sums of degree $n$ are
linearly independent. The degree of $X_T$ is~$n$.
\end{proof}

\noindent The Lean proof does not use power sums. It shows directly that two
trees with the same colouring counts have the same number of ordered partitions
into connected blocks of every sequence of sizes, by counting colourings of
forests ($q^c(q-1)^{n-c}$ proper $q$-colourings for a forest with $c$
components) with refined colour classes (Section~\ref{sec:lean}).

\begin{lemma}[number of leaves; {\cite[Proposition~2(v)]{mmw}}]\label{lem:leaves}
If $n\ge3$, then $N_T((n-1,1))$ is the number of leaves of $T$. Hence $X_T$
determines the number $k=|K(T)|$ of core vertices.
\end{lemma}

\begin{proof}
A partition into a connected block of size $n-1$ and a single vertex $u$ exists
exactly when $T-u$ is connected, that is, when $u$ is a leaf; it is then unique.
\end{proof}

\begin{lemma}[properness and weights; {\cite[Theorem~4.15, Corollary~4.20]{got}}]\label{lem:weights}
Let $T$ and $T'$ be trees with $X_T=X_{T'}$ and $n\ge3$ vertices. If $T$ is
proper, then $T'$ is proper and $W(T')=W(T)$.
\end{lemma}

\begin{proof}
Let $k$ be the common core order (Lemma~\ref{lem:leaves}). In a partition of
$V(T)$ into connected blocks with no block of size~$1$, each block contains a
core vertex (two leaves are never adjacent when $n\ge3$), so there are at most
$k$ blocks. If there are exactly $k$, each block contains exactly one core
vertex~$v$ together with leaves, and each leaf lies in the block of its only
neighbour; so the block of $v$ consists of $v$ and all its leaf neighbours, and
there is no block of size~$1$ only if $T$ is proper. Hence a tree has such a
partition with $k$ blocks if and only if it is proper, and then the partition is
unique and its block sizes are the core weights. By Lemma~\ref{lem:stanley}
these numbers agree for $T$ and~$T'$.
\end{proof}

\noindent In particular there is a weight-preserving bijection between the cores
of $T$ and $T'$; it need not preserve adjacency.

\begin{lemma}[core cuts]\label{lem:cuts}
Let $T$ be proper with $k\ge2$. Removing a core edge $e$ splits the core into
two sides $K_1,K_2$, and $V(T)$ into two connected blocks of sizes
$w(K_1),w(K_2)\ge2$. For trees $T,T'$ as in Lemma~\ref{lem:weights}, the
multisets of unordered pairs $\{w(K_1),w(K_2)\}$ over all core edges of $T$
and of $T'$ coincide.
\end{lemma}

\begin{proof}
Two-block connected partitions of $V(T)$ correspond to the edges of $T$; leaf
edges give the sizes $(n-1,1)$, core edges give two sizes at least~$2$, since
each side contains a core vertex of weight at least~$2$. Apply
Lemma~\ref{lem:stanley} to the partitions $(a,n-a)$ with $2\le a\le n-2$; $T'$
is proper by Lemma~\ref{lem:weights}.
\end{proof}

\section{Proof of Theorem~\ref{thm:A}}\label{sec:A}

Let $T$ satisfy the hypotheses and $X_{T'}=X_T$. By Lemma~\ref{lem:weights}, $T'$
is proper with $W(T')=W(T)$; fix a weight-preserving bijection
$\varphi:K(T)\to K(T')$. Let $A\subseteq K(T)$ be the set of core leaves of $T$
and $A'=\varphi(A)$. By~(A2) each $a\in A$ is the only core vertex of $T$ with
weight $w(a)$, so $\varphi(a)$ is the only core vertex of $T'$ with that weight.

\emph{Decoding.} For an integer $s$ that is the total weight of some
$P\subseteq K(T)$, let $\mathrm{marks}(s)=P\cap A$; by~(A3) this does not depend
on~$P$. If $P'\subseteq K(T')$ has $w(P')=s$, then $\varphi^{-1}(P')$ has total
weight $s$, so $P'\cap A'=\varphi(\mathrm{marks}(s))$. Thus in both trees the
total weight of a set of core vertices determines which anchors it contains.

\emph{Every side contains an anchor.} In $T$, each side of a core edge contains
a leaf of the core (follow a longest path inside the side, starting at the
endpoint of the edge), hence an anchor. If $e'$ is a core edge of $T'$ with side
weights $\{a,b\}$, then by Lemma~\ref{lem:cuts} some core edge of $T$ has the same
side weights, so $\mathrm{marks}(a)$ and $\mathrm{marks}(b)$ are nonempty and
complementary in~$A$; by decoding, the sides of $e'$ contain the anchors
$\varphi(\mathrm{marks}(a))$ and $\varphi(\mathrm{marks}(b))$, both nonempty. No
assumption on the core leaves of $T'$ is needed.

\emph{Records.} Fix $r\in A$ and $r'=\varphi(r)$, and root both cores there. For
a core vertex $v\ne r$ of $T$, let $D_v$ be the set of descendants of $v$ (its
side of the edge to its parent), and put $D_r=K(T)$. The \emph{record} of $v$ is
$(\mathrm{marks}(w(D_v)),\,w(D_v))$; for $T'$ we define records in the same way,
with marks pulled back by~$\varphi$. Each record is determined by the side weights
of one core edge: of the two sides, exactly one contains $r$, and it is the one
whose marks contain~$r$. By Lemma~\ref{lem:cuts}, the multisets of records of the
non-root vertices of $T$ and $T'$ coincide, and both roots have the record
$(A,n)$.

\emph{Order.} The sets $D_v$ form a laminar family: two of them are nested or
disjoint. We claim that $D_u\subsetneq D_v$ holds if and only if either
$\mathrm{marks}(u)\subsetneq\mathrm{marks}(v)$, or the marks are equal and
$w(D_u)<w(D_v)$ (writing $\mathrm{marks}(u)$ for the marks of the record of $u$).
If $D_u\subsetneq D_v$, the marks are included and the weight is smaller, since
weights are positive. Conversely, nonempty marks that are included or equal force
$D_u\cap D_v\ne\emptyset$, so the sets are nested; $D_v\subseteq D_u$ would give
the opposite inclusion of marks or a weight that is not smaller. It follows that
different vertices have different records, so the records identify the core
vertices of $T$, and in the same way those of $T'$; this gives a bijection
$\beta:K(T)\to K(T')$ with equal records and $\beta(r)=r'$, and $\beta$
preserves the relation $D_u\subsetneq D_v$.

\emph{Adjacency and weights.} The parent of $u\ne r$ is the vertex $v$ for which
$D_v$ is the smallest set strictly containing $D_u$, so $\beta$ preserves
parents, and hence adjacency. The weight of $v$ is
$w(D_v)-\sum_u w(D_u)$ over the children $u$ of~$v$, so $\beta$ preserves
weights. Thus $\beta$ is an isomorphism of weighted cores, and attaching $w(v)-1$
leaves at every core vertex gives $T\cong T'$. \qed

\begin{remark}[the Lean proof]\label{rem:leanA}
The Lean proof follows the same route with a slightly different bookkeeping:
it compares the sets of records of both trees, associates to each core vertex
the smallest record containing it, recovers adjacency from the fact that two
core vertices are adjacent exactly when one cut separates them, and recovers
the weights by induction on the size of the branches. The written proof prepared
by Codex allows an arbitrary set of uniquely weighted anchors containing the core
leaves; this is equivalent to~(A2)--(A3), since every condition on the anchors is
universal over the anchors.
\end{remark}

\section{Proof of Theorem~\ref{thm:B}}\label{sec:B}

\paragraph{Coloured connected sets.} A \emph{coloured tree} is a finite tree $R$
with a map $c$ from its vertices to an alphabet. Its \emph{connected-set
polynomial} is
\[
C(R)=\sum_{\emptyset\ne B\subseteq V(R),\ R[B]\text{ connected}}\ \ \prod_{v\in B}z_{c(v)},
\]
in commuting variables $z_\alpha$, one for each colour. A leaf of $R$ is
\emph{tagged} if its colour occurs at no other vertex of~$R$.

\begin{lemma}[extraction]\label{lem:extract}
Let $T$ satisfy (B1) and (B2), and let $X_{T'}=X_T$. Colour each core vertex by
its weight. Then the cores of $T$ and $T'$ have the same connected-set
polynomial.
\end{lemma}

\begin{proof}
By Lemma~\ref{lem:weights}, $T'$ is proper with the same multiset $W$ of core
weights; condition (B2) depends only on $W$, so it also holds for~$T'$. In a
proper tree, a partition of $V(T)$ into connected blocks of sizes at least~$2$
puts each leaf in the block of its neighbour, so it corresponds to a partition of
the core into connected blocks, a block $B$ having size $w(B)$. Let $S$ be the
set of values in~$W$. By~(B2), a core block has total weight in $S$ if and only
if it is a single vertex. Let $M$ be a submultiset of $W$ with $|M|\ge2$ and
$\lambda$ the partition with parts $W\setminus M$ and $\sum M$. The partitions of
$V(T)$ counted by $N_T(\lambda)$ have exactly one block whose size is not in~$S$;
on the core this is one connected set $B$ with $|B|\ge2$, all other core vertices
being singletons whose weights form $W\setminus M$, so the weights in $B$ form the
multiset~$M$. Hence $N_T(\lambda)$ is the number of connected sets of core
vertices whose multiset of weights is~$M$, and the same holds for $T'$ with the
same~$\lambda$. Sets with one vertex are counted by~$W$. By
Lemma~\ref{lem:stanley} all coefficients of $C$ agree.
\end{proof}

\begin{lemma}[peeling]\label{lem:colour}
Let $R$ be a coloured tree with at most one untagged leaf. If $R'$ is any
coloured tree with $C(R')=C(R)$, then there is a colour-preserving isomorphism
$R\to R'$.
\end{lemma}

\begin{proof}
Induction on the number $m$ of vertices of~$R$. The linear part of $C$ lists the
colours with multiplicity, so $R'$ has $m$ vertices and the same colours; for
$m=1$ we are done. Let $m\ge2$. Then $R$ has at least two leaves, so it has a
tagged leaf $a$, of colour~$\alpha$; $R'$ has exactly one vertex $a'$ of
colour~$\alpha$. The connected sets of size~$2$ are the edges, so the sum of the
coefficients of the monomials $z_\alpha z_\gamma$ is the degree of the vertex of
colour $\alpha$; hence $a'$ is a leaf of~$R'$.

Let $v$ be the neighbour of $a$, of colour $\gamma$, and $R_0=R-a$. Setting
$z_\alpha=0$ in $C(R)$ gives $C(R_0)$. Every connected set containing $a$ is
$\{a\}$ or $\{a\}\cup B$ with $B$ connected in $R_0$ and $v\in B$, so
\[
F=\frac{C(R)-C(R_0)}{z_\alpha}-1
\]
is the connected-set polynomial of the sets of $R_0$ that contain~$v$. Its
linear part is $z_\gamma$, so the neighbour $v'$ of $a'$ in $R'$ also has
colour~$\gamma$, and the same formula, applied to $R'$, gives the same~$F$. Now
recolour $v$ with a new colour $y$ that occurs in neither tree; the resulting tree
$R_0^y$ has
\[
C(R_0^y)=C(R_0)-F+z_y\,F/z_\gamma,
\]
since every connected set containing $v$ contributes one factor $z_\gamma$ for
$v$, which is replaced by $z_y$. The same holds for $R'_0{}^{y}$, obtained from
$R'-a'$ by recolouring $v'$, so $C(R_0^y)=C(R'_0{}^{y})$.

In $R_0^y$, a leaf is either a leaf of $R$ other than $a$, whose colour is still
unique if it was unique in $R$, or the vertex $v$, whose colour $y$ is unique.
So $R_0^y$ has at most one untagged leaf, and by induction there is a
colour-preserving isomorphism $R_0^y\to R'_0{}^{y}$. It maps $v$ to $v'$, the
only vertices of colour~$y$; restoring the colour $\gamma$ and mapping $a$ to
$a'$ gives the required isomorphism.
\end{proof}

\begin{proof}[Proof of Theorem~\ref{thm:B}]
By Lemma~\ref{lem:extract} the weight-coloured cores of $T$ and $T'$ have the
same connected-set polynomial. By~(B3) the core of $T$ has at most one untagged
leaf, so by Lemma~\ref{lem:colour} the cores are isomorphic by a
weight-preserving isomorphism. Attaching $w(v)-1$ leaves at each core vertex
gives $T\cong T'$.
\end{proof}

\noindent Lemma~\ref{lem:colour} assumes nothing about $R'$; that is why no
hypothesis on $T'$ is needed. In Lean, the lemma is proved for trees coloured by
natural numbers, so that a new colour always exists.

\section{Examples and estimates (not formalized)}\label{sec:examples}

Nothing in this section is part of the Lean verification. The general
constructions are derived by hand; the specific trees were checked by Python
programs against the literal Lean hypotheses, by the reviewers of the work and
again while writing this paper.

\begin{proposition}[size of the families]\label{prop:size}
Let $T$ have $n$ vertices and $\ell$ core leaves. If $T$ satisfies (A2)--(A3),
then $2^\ell\le n+1$. If $T$ satisfies (B1) and (B3), then $\ell<\sqrt{2n+2}$.
\end{proposition}

\begin{proof}
By (A3), different sets of core leaves have different total weights, which lie
in $\{0,\dots,n\}$. Under (B3), at least $\ell-1$ core leaves have pairwise
distinct weights, all at least~$2$, so $n\ge2+3+\dots+\ell=(\ell-1)(\ell+2)/2$,
which gives $\ell^2+\ell\le 2n+2$.
\end{proof}

\begin{example}[family~(A)]\label{ex:A}
Let $S$ be any finite tree with at least two vertices. Give its non-leaves
arbitrary weights at least~$2$, with total $s$, choose $B>s$ with $B\ge2$, and
give the leaves of $S$ the weights $B,2B,4B,\dots$. Every subset sum of the
internal weights is less than $B$, so the total weight of a set $P$ determines
$\lfloor w(P)/B\rfloor$ and hence its leaves of $S$; and the leaf weights exceed
all internal weights. Attaching $w(v)-1$ leaves gives a tree satisfying
Theorem~\ref{thm:A} with core~$S$. So every tree shape occurs as a core, with any
pattern of internal weights; by Proposition~\ref{prop:size} the number of
vertices grows exponentially in the number of core leaves. The smallest member
found by the final reviewer that is not a caterpillar or a spider and has
diameter at least~$6$ (exhaustive over cores with at most $8$ vertices; the
reviewer of this paper confirmed it over all cores) has $33$ vertices: its core
is a spider with legs of lengths $2,2,1$, the three internal core vertices have
weight~$2$, and the core leaves have weights $3,8,16$. It has diameter~$6$, so
it is covered by the later public results~\cite{kyf} and~\cite[pull
request~1081]{sc}. Among the
members that are not caterpillars, the smallest of diameter at least~$7$ have
$41$ vertices (exhaustive search over the numbers and weights of internal
vertices and core leaves, which is all that (A1)--(A3) depend on). One of them
has as core a path $p_1p_2p_3p_4$ with core leaves $a$ at $p_1$, $b$ at $p_4$ and
$c$ at $p_2$, the weights $2$ at $p_1,\dots,p_4$, and the weights $3,10,20$ at
$a,b,c$; it has diameter~$7$ and seven vertices of degree at least~$3$, and also
satisfies the hypotheses of Theorem~\ref{thm:B}. Remark~\ref{rem:public} compares
it with the public criteria.
\end{example}

\begin{example}[family~(B)]\label{ex:B}
Let the core be the complete binary tree of height $h\ge3$, with $\ell=2^h$
leaves, and put $b=\ell+2$. Give all internal core vertices and one core leaf the
weight $b$, and the other core leaves the weights $b+1,\dots,b+\ell-1$. All
weights lie in $[b,2b)$, so (B2) holds, and only one core leaf has a repeated
weight. The tree has $(2\ell-1)b+\ell(\ell-1)/2$ vertices and diameter $2h+2$,
and its number of core leaves grows like the square root of the number of
vertices. For $h=3$ it has $178$ vertices and diameter~$8$; the weight $10$
occurs eight times, once at a core leaf. It violates~(A2), the internal weights
repeat, and the internal vertices of equal weight are adjacent.
\end{example}

\begin{example}[small members of family~(B)]\label{ex:small}
Take as core the spider with legs of lengths $2,2,1$ (a centre $c$, paths
$c x_1 a_1$ and $c x_2 a_2$, and a leaf $a_3$ at~$c$). With weights $5$ at $c,x_1,x_2$
and $5,4,6$ at $a_1,a_2,a_3$, the tree has $30$ vertices; with $6,6,6$ and
$6,5,7$ it has $36$ vertices. Both satisfy (B1)--(B3), have diameter~$6$, and
violate (A2). The final reviewer of Theorem~\ref{thm:B} found that they lie
outside all published and public criteria it checked, among them those
of~\cite{zeng,awan}, \cite[Corollary~5.15]{got}, diameter at most~$5$,
caterpillars, spiders, \cite{aur} and~\cite{sc} (as of pull request~988
of~\cite{sc}). Both have diameter six, so the later
results~\cite{kyf} and~\cite[pull request~1081]{sc} cover them. Among the members
of family~(B) with at most $20$ vertices, those that are not caterpillars or
spiders and have diameter at least~$6$ all have $18$ vertices (for example the
core above with weights $3,3,3$ and $3,2,4$); they lie outside Zeng's class and
family~(A) and are covered by~\cite{hj}. Among the members of family~(B) outside
Zeng's class that are not caterpillars, the smallest of diameter at least~$7$
have $21$ vertices (covered by~\cite{hj}). Members of diameter~$7$ with $30$
vertices also exist: take as core a path $p_0p_1p_2p_3$ with one further core
vertex attached at each of $p_0$, $p_1$, $p_3$, give $p_0,\dots,p_3$ weight~$3$
and the three core leaves (at $p_0,p_1,p_3$) the weights $3,4,11$. This tree
satisfies (B1)--(B3), violates~(A2), lies outside Zeng's class, and has seven
vertices of degree at least~$3$; Remark~\ref{rem:public} compares it with the
public criteria.
\end{example}

\begin{example}[the families are incomparable]\label{ex:incomp}
A path core with weights $16,2,4,2,32$ (a caterpillar with $56$ vertices)
satisfies (A1)--(A3) but violates (B2), since $4=2+2$. The tree with $178$
vertices of Example~\ref{ex:B} satisfies (B1)--(B3) but violates~(A2).
\end{example}

\begin{remark}[comparison with public criteria]\label{rem:public}
While revising this paper on 27 November 2026 (after 18:42~UTC), Claude Code
checked the trees with $41$ vertices (Example~\ref{ex:A}) and $30$ vertices and
diameter~$7$ (Example~\ref{ex:small}), with a program of its own, against the
explicit criteria of the works cited here, with the following results. Both
trees are proper, are not caterpillars or spiders, have diameter~$7$ and more
than $29$ vertices, and have a repeated internal weight; so they are outside
\cite{mmw,apz,ls,hj}, Zeng's and Awan's classes, \cite[Corollary~5.15]{got},
the diameter-five results of~\cite{got,apdmoz}, and the diameter-six results~\cite{kyf} (checked on 28 November)
and~\cite[pull request~1081]{sc}; each has seven vertices of degree at least~$3$,
so~\cite{carp} (checked on 28 November) does not apply. Each has a unique centroid with at most two
leaf neighbours, whereas its largest branch has $20$ and $14$ vertices,
respectively, so the leaf-rich centroid criterion of~\cite[pull request~977]{sc}
fails; no edge of either tree satisfies the balanced-edge criterion
of~\cite[pull request~977]{sc} or the subdivided-edge criteria of~\cite[pull
requests~972 and~977]{sc}; no vertex has all its branches isomorphic, as the
repeated-branch criteria of~\cite[pull request~972]{sc} require; and every vertex
of degree at least~$3$ has a leaf neighbour, so neither tree has the long arms at
its branch vertices required in~\cite[pull requests~985, 1024 and~1029]{sc}. The
observable path-count certificate of~\cite[pull request~1030]{sc} fails for
both. Both fail necessary conditions of every criterion of~\cite{aur} listed in
the prior-work check of the final reviewer (repeated weights at adjacent vertices,
with multiplicity at least four, and a vertex of repeated weight with no neighbour
of unique weight). The density and genericity results of~\cite{sc}, for random or generic edge lengths, do not decide individual trees
and were not checked. On 28 November 2026 the release-day search (see the end of
the paper) added Atwell's criteria: both trees violate the condition
of~\cite[Theorem~2.1]{atwell1} that every peripheral mass be smaller than the
repeated weight (the repeated weight is $2$ with peripheral masses $3,10,20$ in the tree with $41$
vertices, and $3$ with peripheral masses $4,11$ in the tree with $30$ vertices), and both
violate the hypotheses of~\cite{atwell2} (vertices of repeated weight are adjacent,
and the weights are not all congruent modulo any $m\ge2$); the tree with $178$
vertices of Example~\ref{ex:B} also violates the former (repeated weight $10$,
peripheral masses $11,\dots,17$).
\end{remark}

\begin{remark}[computer checks]\label{rem:brute}
The final reviewers also ran exhaustive checks (checks, not proofs): for every
$n$ from $4$ to $18$, no member of family~(A) shares $X$ with another tree; for
$n\le20$, no member of family~(B) shares $X$ with another proper tree with the
same core weights; and Lemma~\ref{lem:colour} was tested on all coloured trees
with at most $8$ vertices and at most $3$ colours, where collisions of $C$ occur
only among trees with at least two untagged leaves.
\end{remark}

{\sloppy
\section{Formal verification}\label{sec:lean}

Theorems~\ref{thm:A} and~\ref{thm:B} are formally verified in Lean~4 (version
4.34.1) with Mathlib~\cite{mathlib} (commit \texttt{d13f23b7}), as the
statements \texttt{AnchorReconstruction} and \texttt{PeelingReconstruction} in
the namespace \texttt{StanleyContinuation}. The $38$ files ($3{,}222$ lines) are
in version~\repoversion{} of the repository~\cite{repo}, directory
\texttt{lean/Research/TaskV/Stanley/}; records of the freezing and of the checks
are in the folder \texttt{certificates/stanley-chromatic-trees}. The import
closure of Theorem~\ref{thm:A} has $26$ files, that of Theorem~\ref{thm:B} has
$30$, and $18$ files are shared.

\paragraph{The statements.} The statement files \texttt{ChallengeA.lean} and
\texttt{ChallengePeel.lean}, in ASCII transliteration (the \texttt{import},
\texttt{namespace}, \texttt{noncomputable section}, \texttt{open} and
\texttt{end} lines, comments and blank lines are omitted; the files use Unicode:
\texttt{Nat} stands for $\N$, \texttt{->} for $\to$, \texttt{forall},
\texttt{exists}, \texttt{in}, \texttt{/\char`\\}, \texttt{<->}, \texttt{<=},
\texttt{!=}, \texttt{subset}, \texttt{inter} and \texttt{sum} for $\forall$,
$\exists$, $\in$, $\wedge$, $\leftrightarrow$, $\le$, $\ne$, $\subseteq$, $\cap$
and $\sum$, \texttt{\char`\\} is set difference, and \texttt{G =\char`~g H}
stands for Mathlib's type $G\simeq_g H$ of graph isomorphisms):
{\footnotesize
\begin{verbatim}
-- ChallengeA.lean
def colorCoefficient {n : Nat} (G : SimpleGraph (Fin n)) (k : Nat)
    (multiplicity : Fin k -> Nat) : Nat :=
  Nat.card { f : Fin n -> Fin k //
    (forall i j, G.Adj i j -> f i != f j) /\
    forall c, Nat.card {i : Fin n // f i = c} = multiplicity c }
def SameCSF {n m : Nat} (G : SimpleGraph (Fin n)) (H : SimpleGraph (Fin m)) : Prop :=
  forall k (multiplicity : Fin k -> Nat),
    colorCoefficient G k multiplicity = colorCoefficient H k multiplicity
def neighbors {n : Nat} (G : SimpleGraph (Fin n)) (v : Fin n) : Finset (Fin n) := by
  classical exact univ.filter (G.Adj v)
def core {n : Nat} (G : SimpleGraph (Fin n)) : Finset (Fin n) := by
  classical exact univ.filter fun v => (neighbors G v).card != 1
def coreWeight {n : Nat} (G : SimpleGraph (Fin n)) (v : Fin n) : Nat :=
  1 + ((neighbors G v) \ core G).card
def DecodableLeafAnchors {n : Nat} (G : SimpleGraph (Fin n)) : Prop :=
  2 <= (core G).card /\
  (forall v in core G, 2 <= coreWeight G v) /\
  exists A : Finset (Fin n),
    A subset core G /\
    (forall a in A, forall v in core G, coreWeight G v = coreWeight G a -> v = a) /\
    (forall v in core G, ((neighbors G v) inter core G).card = 1 -> v in A) /\
    forall P Q : Finset (Fin n), P subset core G -> Q subset core G ->
      (sum v in P, coreWeight G v) = (sum v in Q, coreWeight G v) ->
      forall a in A, (a in P <-> a in Q)
def AnchorReconstruction : Prop :=
  forall (n m : Nat) (G : SimpleGraph (Fin n)) (H : SimpleGraph (Fin m)),
    G.IsTree -> H.IsTree -> DecodableLeafAnchors G -> SameCSF G H ->
      Nonempty (G =~g H)
-- ChallengePeel.lean (imports ChallengeA.lean)
def untaggedCoreLeaf {n : Nat} (G : SimpleGraph (Fin n)) (v : Fin n) : Prop :=
  v in core G /\ ((neighbors G v) inter core G).card = 1 /\
    exists u in core G, u != v /\ coreWeight G u = coreWeight G v
def PeelingWeights {n : Nat} (G : SimpleGraph (Fin n)) : Prop :=
  2 <= (core G).card /\
  (forall v in core G, 2 <= coreWeight G v) /\
  (forall v in core G, forall P : Finset (Fin n), P subset core G -> 2 <= P.card ->
    (sum u in P, coreWeight G u) != coreWeight G v) /\
  forall v u, untaggedCoreLeaf G v -> untaggedCoreLeaf G u -> v = u
def PeelingReconstruction : Prop :=
  forall (n m : Nat) (G : SimpleGraph (Fin n)) (H : SimpleGraph (Fin m)),
    G.IsTree -> H.IsTree -> PeelingWeights G -> SameCSF G H ->
      Nonempty (G =~g H)
\end{verbatim}}

\noindent \texttt{colorCoefficient G k mult} counts the proper colourings
$f:V\to\{1,\dots,k\}$ with exactly $\mathrm{mult}(c)$ vertices of colour~$c$,
which is the coefficient of $x_1^{\mathrm{mult}(1)}\cdots x_k^{\mathrm{mult}(k)}$
in $X_G$ (a proper colouring with these class sizes uses only the colours
$1,\dots,k$). Since $X_G$ is symmetric and each of its monomials involves finitely
many variables, \texttt{SameCSF G H}, which compares these numbers for all $k$
and all sequences (zeros included), is equivalent to $X_G=X_H$; it also forces
equal numbers of vertices (Lean: \texttt{sameOrder\_of\_sameCSF}). The final
reviewer of Theorem~\ref{thm:A} confirmed the equivalence numerically on all
$52$ graphs with at most $5$ vertices, including Stanley's pair of
non-isomorphic graphs with the same~$X$. Every finite tree is isomorphic to a
tree on \texttt{Fin n}, and Mathlib's \texttt{IsTree} means connected and
acyclic. \texttt{core G} is the set of vertices whose degree is not~$1$; for
trees with at least three vertices these are the non-leaves, and trees with at
most two vertices have at most one such vertex, so they never satisfy the
hypotheses. \texttt{coreWeight} is $1$ plus the number of neighbours outside the
core, that is, the weight~$w$. In \texttt{DecodableLeafAnchors} the set
\texttt{A} must contain every core vertex with exactly one core neighbour, and
all conditions on \texttt{A} are universal over its elements; so the predicate is
equivalent to (A1)--(A3) with $A$ the set of core leaves (the final reviewer also
confirmed this by exhaustive search for $n\le10$). \texttt{untaggedCoreLeaf}
expresses a core leaf whose weight occurs at another core vertex, and
\texttt{PeelingWeights} is (B1)--(B3). The hypotheses constrain \texttt{G} only,
\texttt{H} is an arbitrary tree, and the conclusion is an isomorphism of the
whole trees, leaves included.

\paragraph{Review and freezing.} The statement files were read against the
written proofs by Codex agents that had not written them (approvals at
03:47~UTC for \texttt{ChallengeA.lean} and 03:51~UTC for
\texttt{ChallengePeel.lean} on 21 November 2026), compiled with no diagnostics
(03:53~UTC), and frozen at 03:53:48 and 03:53:59~UTC by recording their SHA-256
hashes (\texttt{683c4463\ldots} and \texttt{dcc0d30a\ldots}) together with those
of the written proofs, the reviews and the compilation receipts. According to the
file timestamps, all proof files were created after the freezes (the earliest at
03:54:26~UTC), with one exception: \texttt{ProfileDefs.lean}, an $18$-line file
of two definitions (counts of ordered partitions into connected blocks), was
created at 03:48:33~UTC and not changed afterwards; the statement files do not
import it. The statement files were not changed after the freezes; the final
reviewers recomputed their hashes.

\paragraph{The proofs.} The check files \texttt{CheckAnchor.lean} and
\texttt{CheckPeel.lean} contain, apart from the \texttt{import} and
\texttt{namespace} lines, exactly
{\footnotesize
\begin{verbatim}
theorem check_anchorReconstruction : AnchorReconstruction := anchorReconstruction
theorem check_peelingReconstruction : PeelingReconstruction := peelingReconstruction
\end{verbatim}}
\noindent respectively. Table~\ref{tab:lean} lists the Lean counterparts of the
steps (names in the namespace \texttt{StanleyContinuation}; file names in
parentheses, without \texttt{.lean}).

\begin{table}[!htbp]
\centering\footnotesize
\renewcommand{\arraystretch}{1.05}
\begin{tabular}{@{}>{\raggedright\arraybackslash}p{0.27\textwidth}>{\raggedright\arraybackslash}p{0.69\textwidth}@{}}
\hline
Step of the written proof & Lean (file) \\
\hline
Lemma~\ref{lem:stanley} & \texttt{sameConnected\_of\_sameCSF} (\texttt{CSFBridge}),
with \texttt{Forest}, \texttt{Isolation}, \texttt{Cell}, \texttt{Refinement};
\texttt{sameOrder\_of\_sameCSF} (\texttt{Histogram}) \\
Lemma~\ref{lem:leaves} & \texttt{singleton\_count\_eq\_leaf\_count},
\texttt{same\_core\_card} (\texttt{AnchorLeaves}); \texttt{leafCount\_of\_sameCSF}
(\texttt{CoreWeights}) \\
Lemma~\ref{lem:weights} & \texttt{coreWeights\_of\_sameCSF} (\texttt{CoreWeights}),
\texttt{weightPreservingCoreEquiv} (\texttt{StarProfiles}) \\
Lemma~\ref{lem:cuts} and decoding & \texttt{matching\_core\_cut} (\texttt{StarSplit}),
\texttt{anchorCutProfiles\_of\_sameCSF} (\texttt{AnchorCutLabels}) \\
Records and order & \texttt{matching\_record\_order} (\texttt{AnchorOrder}), with
\texttt{AnchorTerminals}, \texttt{AnchorSides}, \texttt{AnchorTree},
\texttt{AnchorIncidence} \\
Adjacency and weights & \texttt{reconstruct\_core\_adj} (\texttt{AnchorRecords}),
\texttt{weights\_of\_record\_sums} (\texttt{AnchorVertexWeights}) \\
Attaching the leaves & \texttt{expand\_weighted\_core\_iso} (\texttt{ExpandCore}) \\
Theorem~\ref{thm:A} & \texttt{anchorReconstruction} (\texttt{AnchorProof}),
\texttt{check\_anchorReconstruction} (\texttt{CheckAnchor}) \\
\hline
Lemma~\ref{lem:extract} & \texttt{coreColoredSets\_of\_sameCSF},
\texttt{peel\_coloredSets\_of\_sameCSF} (\texttt{PeelInventory}), with
\texttt{PeelPartitions}, \texttt{PeelLabelings} \\
Lemma~\ref{lem:colour}, one step & \texttt{peel\_step},
\texttt{core\_peel\_step\_of\_sameCSF} (\texttt{PeelRetag}), with \texttt{PeelTagged},
\texttt{PeelDeletion} \\
Lemma~\ref{lem:colour}, induction & \texttt{atMostOne\_after\_deletion}
(\texttt{PeelLeaves}), \texttt{reattach\_tagged\_leaf} (\texttt{PeelReattach}),
\texttt{colored\_tree\_reconstruction} (\texttt{PeelInduction}) \\
Theorem~\ref{thm:B} & \texttt{peelingReconstruction} (\texttt{PeelProof}),
\texttt{check\_peelingReconstruction} (\texttt{CheckPeel}) \\
\hline
\end{tabular}
\caption{Lean counterparts of the steps of the proofs.}
\label{tab:lean}
\end{table}

\noindent In \texttt{PeelInventory}, the multiplicities of repeated weights are
handled by naming the vertices outside the connected set; each set is then
counted as many times as there are weight-preserving permutations of these
names, a factor that is the same for both trees and is cancelled. The axioms of
\texttt{check\_anchorReconstruction} and \texttt{check\_peelingReconstruction}
are exactly \texttt{propext}, \texttt{Classical.choice} and \texttt{Quot.sound}.
No file uses \texttt{sorry}, \texttt{native\_\allowbreak decide},
\texttt{bv\_\allowbreak decide}, \texttt{decide}, an \texttt{axiom} declaration or
\texttt{set\_option}, and the checked package sets no options.

\paragraph{Checks.} Each theorem was checked in two new empty directories on the
author's Mac, one after the other, with the pinned toolchain, on 27 November
2026: each directory compiled every module of the import closure separately
with warnings treated as errors (\texttt{lake build --wfail}), built the
package's root module \texttt{Stanley}, which imports only the check file (each
directory contained exactly the import closure),
replayed every module with \texttt{leanchecker}, and audited the axioms of all public theorems and lemmas ($86$ for
Theorem~\ref{thm:A}, $92$ for Theorem~\ref{thm:B}). For Theorem~\ref{thm:A} ($27$ modules, $57$ commands per
directory) the directories ran from 03:01:56 to 03:09:56 and from 03:09:56 to
03:17:37~UTC; for Theorem~\ref{thm:B} ($31$ modules, $65$ commands) from 05:11:58
to 05:20:33 and from 05:20:33 to 05:27:37~UTC. All commands ended with exit
code~$0$, all \texttt{leanchecker} runs printed nothing, the two directories
produced byte-identical \texttt{.olean} files, every audited theorem uses
only the three axioms above, and the axiom lists agreed. Intermediate results had been compiled and replayed with \texttt{leanchecker}
during development (21 and 27 November), each round in a new directory that
reused earlier checked modules by hash; these receipts are kept on record. Each final reviewer (see the
next section) checked with plain \texttt{lean}, against the compiled files of the earlier of the two directories and without writing to the project, that the check theorem
has exactly the frozen type and the three axioms, and that the definitions are
unchanged; the reviewers did not recompile the files. Mathlib and its
dependencies were taken from a shared prebuilt cache at the pinned commits and
were not rechecked; the \texttt{leanchecker} replays cover the project modules.
In the repository the files have module names under
\texttt{Research.TaskV.Stanley}, keep the namespace \texttt{StanleyContinuation},
and differ from the checked files only in their \texttt{import} lines
(\texttt{ChallengeA.lean} and \texttt{Isolation.lean}, which import only Mathlib,
are byte-identical); the root module \texttt{Stanley} of the checked package is
not part of the repository.
These versions were compiled and replayed with \texttt{leanchecker} again, module by module, in two new empty directories on the author's Mac (27 November 2026, 18:11--18:18~UTC), with no warnings and identical \texttt{.olean} files; the axioms of the two check theorems are again exactly \texttt{propext}, \texttt{Classical.choice} and \texttt{Quot.sound}.
In the whole-project rebuild of version~\repoversion{} of the repository (on the author's Windows computer, in one pass from 27 to 28 November 2026), these files were compiled and replayed with \texttt{leanchecker} again together with the rest of the project; no module needed a second attempt (details in the repository). The formalization has not been checked with comparator or with an
independent kernel.

\par}

\section*{How this paper was produced}
This work was carried out on 21 November and from 27 to 28 November 2026 (UTC), and the paper was written on 27 November and revised on 28 November, in a workflow directed by the author, using OpenAI Codex
with the model GPT-6.2-Astra (reasoning level ultra) and Anthropic Claude Code
with the model Claude Opus~5.5. The author asked for a new search for problems,
with a higher threshold aimed at famous open problems, and Claude Code passed
this request to Codex at 03:19~UTC on 21 November (file time). Codex screened
$94$ problems, wrote notes on eight, and made short attempts on three, among
them Stanley's question, which its report recommended for further work; none of
the three was solved (report at
03:35~UTC). After reading the report, the author asked that Codex continue on
Stanley's question, and Claude Code passed this request to the same Codex session
at 03:43~UTC (the time stamp of its instruction in the session log). According to
the file times, Codex had written the proofs of Theorems~\ref{thm:A}
and~\ref{thm:B} by 03:46 and 03:47~UTC; Codex agents of the same session (model GPT-6.2-Astra) that had not
derived them approved the written proofs at 03:47 and 03:49~UTC, and the
statements were reviewed and frozen as described in Section~\ref{sec:lean}. The
Lean work started at 03:54~UTC; the session reached its usage limit at
04:04~UTC. From then on Codex worked in a queue, one task at a time, and Claude
Code resumed the session at 04:11 and 17:13~UTC on 21 November and at 00:33,
02:21, 02:57, 03:59, 04:36 and 05:09~UTC on 27 November (the time stamps of its
instructions in the session logs); each round ended at a recorded intermediate
result. The Lean proof that equal colouring counts give equal counts of connected
partitions was completed on 21 November (17:13--17:18~UTC), the recovery of the
leaves and of the core weights and the cut records on 27 November
(00:33--02:26~UTC), Theorem~\ref{thm:A} at 03:18~UTC and Theorem~\ref{thm:B} at
05:28~UTC, after the checks of Section~\ref{sec:lean}. The Lean proofs follow the
written proofs (Remark~\ref{rem:leanA}). Codex also ran computations aimed at a
counterexample to the general question, among them a cross-comparison of a
prefix of $1{,}016{,}250{,}000$ of the $14{,}830{,}871{,}802$ trees with $30$
vertices; they found none, cover only part of the trees, and are not part of this
paper. The third, centroid condition mentioned in Section~\ref{sec:intro} and
further unreviewed written work by Codex are not claimed here.

The final reviews were done by Claude Code agents (model Claude Opus~5.5) that
had written none of the files. For Theorem~\ref{thm:A}, Claude Code started a
final review and a separate search for prior work at 03:56~UTC on 27 November;
the search report is dated 04:05~UTC and the review 04:07~UTC. For
Theorem~\ref{thm:B}, one agent did both, from 05:46 to 06:01~UTC. The reviewers
read the statements line by line, read the written proofs and the proof files,
recomputed the hashes, checked the receipts, checked the Lean types and axioms as
described in Section~\ref{sec:lean}, tested the statements and the examples with
programs of their own (Remark~\ref{rem:brute}), and determined the size of the
classes and the smallest examples outside the known criteria. Both decisions were
PASS, with no blocking remarks; their recommended corrections concern the written
account and are taken into account here (the coverage up to $29$ vertices and the
smallest new examples, the bounds on the number of core leaves, the credit for
the known ingredients and the comparison with~\cite{zeng,awan,got}, the
incomparability of the two families, and that the example families are not
verified in Lean). Claude Code (Claude Opus~5.5) checked the sources, checked
the examples of Section~\ref{sec:examples} in Python against the literal Lean
hypotheses, and wrote this text from Codex's manuscripts, the reviews and the
Lean files. An independent Claude Code agent (model Claude Opus~5.5) that had written none
of the files reviewed this paper, and two others reviewed the release (the second
only the parts written after the rebuild of the whole project); their corrections
were applied. Throughout, ``we'' refers to this workflow; the reading
and searching described in the paper were done by the agents. Codex agents
searched for prior work on 21 November (from 03:44~UTC: the arXiv, citation
indexes, GitHub, repositories of AI-assisted mathematics and the public
discussions of the question), repeated a fixed set of $17$ queries daily from 22
to 27 November, and searched again before each round on 27 November; the final
reviewers searched again on 27 November (the arXiv, Semantic Scholar, scite, web
search, GitHub and the public AI projects~\cite{aur,sc}; some code searches
returned nothing because repositories were not indexed, which was not taken as
evidence). On 27 November 2026 (18:42~UTC), while revising this paper after an
independent review, Claude Code read the list of pull requests merged
in~\cite{sc} that day and the documents of those on Stanley's question, and
compared the examples with them (Remark~\ref{rem:public}). Up to 27~November 2026
these searches found no earlier statement of Theorem~\ref{thm:A} or
Theorem~\ref{thm:B}; the related public work found is described in
Section~\ref{sec:intro} (details in the repository).
On 28~November 2026 (08:15--08:34~UTC) Claude Code searched again: the pull
requests, commits and issue of~\cite{sc} (no merge on Stanley's question after
pull request~1081), the repositories of~\cite{aur}, the arXiv (listings up to
26 November; submissions after about 25 November 19:00~UTC were not yet
announced), the OpenAlex and Semantic Scholar lists of works citing the papers of
Martin, Morin and Wagner, Gonzalez, Orellana and Tomba, Zeng and Awan (Semantic
Scholar refused the query for Gonzalez, Orellana and Tomba because of a rate
limit, and an OpenAlex date filter was refused), GitHub, the records of several
AI projects, Zenodo and DataCite. It found no statement of Theorem~\ref{thm:A}
or Theorem~\ref{thm:B}, and four related public works not described before: Atwell's two
preprints~\cite{atwell1,atwell2}, the preprint~\cite{carp} and the
repository~\cite{kyf}, which
are discussed in Section~\ref{sec:intro} and Remark~\ref{rem:public}; none of
them covers the examples of diameter~$7$ given there. This reports the coverage of the searches, not a guarantee
that they were complete. The author is not a professional mathematician.
The author has read this paper and, using a side-by-side table, compared its Lean statements with Theorems~\ref{thm:A} and~\ref{thm:B}. No human mathematician has read the proofs or reviewed this paper.
The correctness of Theorems~\ref{thm:A} and~\ref{thm:B} rests on the Lean proofs
and on the agreement of the Lean statements with the definition of the chromatic
symmetric function; the written proofs are given for the reader. AI systems are
not authors of this paper; the author takes full responsibility for its content.

\medskip
\noindent Thanks are due to Richard Stanley for the question, and to the authors
of~\cite{got,apdmoz,ot,zeng,awan,mmw,hj}, on whose results the comparison rests.

{\small
\begin{thebibliography}{99}
\setlength{\itemsep}{1pt}
\bibitem{apdmoz} J.~Aliste-Prieto, A.~de~Mier, R.~Orellana and J.~Zamora, Marked
graphs and the chromatic symmetric function, \emph{SIAM J. Discrete Math.} 37
(2023), no.~3, 1881--1919, \url{https://doi.org/10.1137/22M148046X};
arXiv:2202.11787v1 (numbering of results as in the arXiv version).
\bibitem{apz} J.~Aliste-Prieto and J.~Zamora, Proper caterpillars are
distinguished by their chromatic symmetric function, \emph{Discrete Math.}
315--316 (2014), 158--164, \url{https://doi.org/10.1016/j.disc.2013.10.016}.
\bibitem{aur} Aurantium-47, First-derivative reconstruction for proper trees with
independent repeated weights, pull request~5 of the repository
\emph{aurantium-mathematics},
\url{https://github.com/Aurantium-47/aurantium-mathematics/pull/5} (working notes,
last updated 22 November 2026).
\bibitem{atwell1} C.~Atwell, Chromatic symmetric reconstruction of proper trees
with leaf-observed repeated-weight cores, research preprint, version 1.0.2,
Zenodo, 12 November 2026, \url{https://doi.org/10.5281/zenodo.24163711}.
\bibitem{atwell2} C.~Atwell, Congruence-class Gram reconstruction for proper
trees: multiple repeated-weight classes without a common-divisor hypothesis,
research preprint, version 1.3.1, Zenodo, 24 October 2026,
\url{https://doi.org/10.5281/zenodo.23627746}.
\bibitem{awan} S.~A. Awan, First-derivative chromatic symmetric reconstruction for
proper trees, arXiv:2609.13464v2 (2026).
\bibitem{buneman} P.~Buneman, The recovery of trees from measures of
dissimilarity, in \emph{Mathematics in the Archaeological and Historical
Sciences}, Edinburgh University Press, 1971, 387--395.
\bibitem{carp} Carptopus, Trees with three branch vertices are determined by
their chromatic symmetric functions, preprint, version 0.1-beta, Zenodo,
17 September 2026, \url{https://doi.org/10.5281/zenodo.22813491}.
\bibitem{repo} H.~Dong, \emph{ai4math-results}, repository,
\url{https://github.com/Horace-Maxwell/ai4math-results}, version~\repoversion{}
(2026); archived at Zenodo under the concept DOI
\url{https://doi.org/10.5281/zenodo.22970254}.
\bibitem{kyf} K.~Feng, Reconstructing trees of diameter six from the ordinary
chromatic symmetric function, manuscript and Lean~4 proof, repository
\url{https://github.com/Kyf-eng/diameter-six-chromatic-trees}, releases v1.0.0
and v1.1.0 (27 November 2026).
\bibitem{got} M.~Gonzalez, R.~Orellana and M.~Tomba, The chromatic symmetric
function in the star-basis, \emph{Discrete Math.} 349 (2026), no.~2, 114691,
\url{https://doi.org/10.1016/j.disc.2025.114691}; arXiv:2404.06002v2 (numbering
of results as in the arXiv version).
\bibitem{hj} S.~Heil and C.~Ji, On an algorithm for comparing the chromatic
symmetric functions of trees, \emph{Australas. J. Combin.} 75 (2019), no.~2,
210--222, \url{https://ajc.maths.uq.edu.au/pdf/75/ajc_v75_p210.pdf}.
\bibitem{ls} M.~Loebl and J.-S. Sereni, Isomorphism of weighted trees and
Stanley's isomorphism conjecture for caterpillars, \emph{Ann. Inst. Henri
Poincar\'e D} 6 (2019), no.~3, 357--384, \url{https://doi.org/10.4171/AIHPD/74};
arXiv:1405.4132v4 (numbering of results as in the arXiv version).
\bibitem{mmw} J.~L. Martin, M.~Morin and J.~D. Wagner, On distinguishing trees by
their chromatic symmetric functions, \emph{J. Combin. Theory Ser.~A} 115 (2008),
no.~2, 237--253, \url{https://doi.org/10.1016/j.jcta.2007.05.008};
arXiv:math/0609339v3 (numbering of results as in the arXiv version).
\bibitem{mathlib}
The mathlib Community, The Lean mathematical library, in \emph{Proc.\ 9th ACM
SIGPLAN Int.\ Conf.\ Certified Programs and Proofs} (CPP '20), ACM, 2020,
367--381, \url{https://doi.org/10.1145/3372885.3373824}.
\bibitem{ot} R.~Orellana and M.~Tomba, Proper implies improper for the chromatic
symmetric function, arXiv:2611.18940v1 (2026).
\bibitem{sc} \emph{solved-conjectures}, repository of AI-assisted mathematics,
\url{https://github.com/1337isnot1337/solved-conjectures}, pull requests 972,
977, 985, 1024, 1029, 1030, 1079 and 1081 on Stanley's question, merged on 27
November 2026 (972 at 04:11, 977 at 04:28, 985 at 05:00, 1024 at 14:31, 1029 at
14:43, 1030 at 14:48, 1079 at 18:00 and 1081 at 18:17~UTC; main branch at commit
\texttt{23c9001e}); written proofs with internal reviews, not peer reviewed and
not formally verified.
\bibitem{stanley} R.~P. Stanley, A symmetric function generalization of the
chromatic polynomial of a graph, \emph{Adv. Math.} 111 (1995), no.~1, 166--194,
\url{https://doi.org/10.1006/aima.1995.1020}.
\bibitem{zeng} Z.~Zeng, Multiplicity-isolated cores and chromatic symmetric
reconstruction of trees, arXiv:2608.18179v1 (2026).
\end{thebibliography}}
\end{document}
